% paper88-copilot-evidence-channels.tex — The Copilot Evidence Channel:
%   Trust-Calibrated LLM Contributions
%   Paper 88 of the Judgment Geometry series.
% Compile: pdflatex paper88-copilot-evidence-channels
\documentclass[11pt]{article}
\usepackage{lmodern}
% jugeo-common.tex — shared preamble for all Judgment Geometry papers
% Include via: % jugeo-common.tex — shared preamble for all Judgment Geometry papers
% Include via: % jugeo-common.tex — shared preamble for all Judgment Geometry papers
% Include via: \input{jugeo-common}

% ── Packages ────────────────────────────────────────────────────────────
\usepackage{amsmath,amssymb,amsthm,mathtools}
\usepackage{tikz-cd}
\usepackage{listings}
\usepackage{xcolor}
\usepackage{xspace}
\usepackage{booktabs}
\usepackage{multirow}
\usepackage{enumitem}
\usepackage[expansion=false]{microtype}
\usepackage{hyperref}
\usepackage{cleveref}
% \usepackage{thmtools}  % disabled: not available in this TeX installation
\usepackage{float}
\usepackage{caption}
\usepackage{subcaption}
\usepackage{tabularx}
\usepackage{stmaryrd}       % \llbracket, \rrbracket
\usepackage{bbm}            % \mathbbm{1}

% ── Page geometry ───────────────────────────────────────────────────────
\usepackage[margin=1in]{geometry}

% ── Theorem environments ────────────────────────────────────────────────
\theoremstyle{plain}
\newtheorem{theorem}{Theorem}[section]
\newtheorem{lemma}[theorem]{Lemma}
\newtheorem{proposition}[theorem]{Proposition}
\newtheorem{corollary}[theorem]{Corollary}

\theoremstyle{definition}
\newtheorem{definition}[theorem]{Definition}
\newtheorem{example}[theorem]{Example}
\newtheorem{construction}[theorem]{Construction}

\theoremstyle{remark}
\newtheorem{remark}[theorem]{Remark}
\newtheorem{notation}[theorem]{Notation}

% ── Python listings style ──────────────────────────────────────────────
\definecolor{jugeoBlue}{HTML}{2563EB}
\definecolor{jugeoGreen}{HTML}{059669}
\definecolor{jugeoRed}{HTML}{DC2626}
\definecolor{jugeoPurple}{HTML}{7C3AED}
\definecolor{jugeoGray}{HTML}{6B7280}
\definecolor{jugeoBg}{HTML}{F8FAFC}

\lstdefinestyle{jugeo-python}{
  language=Python,
  basicstyle=\ttfamily\small,
  keywordstyle=\color{jugeoBlue}\bfseries,
  stringstyle=\color{jugeoGreen},
  commentstyle=\color{jugeoGray}\itshape,
  numberstyle=\tiny\color{jugeoGray},
  backgroundcolor=\color{jugeoBg},
  numbers=left,
  numbersep=6pt,
  frame=single,
  framerule=0.4pt,
  rulecolor=\color{jugeoGray!40},
  breaklines=true,
  breakatwhitespace=true,
  showstringspaces=false,
  tabsize=4,
  xleftmargin=1.5em,
  framexleftmargin=1.5em,
  morekeywords={dataclass,frozen,field,Enum,IntEnum,auto,Any,
                Mapping,Sequence,Optional,match,case},
  literate={→}{{$\to$}}1 {←}{{$\leftarrow$}}1
           {≤}{{$\leq$}}1 {≥}{{$\geq$}}1 {≠}{{$\neq$}}1
           {∈}{{$\in$}}1 {∀}{{$\forall$}}1 {∃}{{$\exists$}}1
           {⊕}{{$\oplus$}}1 {⊖}{{$\ominus$}}1,
  escapeinside={(*@}{@*)},
}
\lstset{style=jugeo-python}

\lstdefinestyle{jugeo-output}{
  basicstyle=\ttfamily\small,
  backgroundcolor=\color{jugeoBg},
  frame=single,
  framerule=0.4pt,
  rulecolor=\color{jugeoGray!40},
  breaklines=true,
  showstringspaces=false,
  xleftmargin=1.5em,
  framexleftmargin=0.5em,
  numbers=none,
}

% ── JuGeo-specific macros ──────────────────────────────────────────────

% Judgment 8-tuple
\newcommand{\J}{\mathcal{J}}
\newcommand{\Jud}[8]{(#1,\,#2,\,#3,\,#4,\,#5,\,#6,\,#7,\,#8)}
\newcommand{\JudShort}[1]{\J_{#1}}

% Components
\newcommand{\coord}{\mathsf{c}}            % coordinate
\newcommand{\prop}{\varphi}                 % proposition
\newcommand{\carrier}{\mathcal{A}}          % carrier type
\newcommand{\evid}{\mathcal{E}}             % evidence bundle
\newcommand{\oblig}{\mathcal{O}}            % obligations
\newcommand{\obstr}{\mathcal{B}}            % obstructions
\newcommand{\trust}{\mathcal{T}}            % trust annotation
\newcommand{\prov}{\Pi}                     % provenance

% Trust algebra
\newcommand{\Talg}{\mathfrak{T}}            % trust ordered algebra
\newcommand{\Eadm}{\mathcal{E}_{\mathrm{adm}}}
\newcommand{\tleq}{\preceq}                 % trust ordering
\newcommand{\tjoin}{\oplus}                 % conservative join
\newcommand{\tatten}{\ominus}               % attenuation
\newcommand{\tpromote}{\uparrow_\pi}        % promotion
\newcommand{\tdemote}{\downarrow_\chi}       % demotion

% Trust tiers
\newcommand{\Tcontra}{\ensuremath{\mathsf{CONTRADICTED}}}
\newcommand{\Tunver}{\ensuremath{\mathsf{UNVERIFIED}}}
\newcommand{\Tcopilot}{\ensuremath{\mathsf{COPILOT}}}
\newcommand{\Toracle}{\ensuremath{\mathsf{ORACLE}}}
\newcommand{\Truntime}{\ensuremath{\mathsf{RUNTIME}}}
\newcommand{\Tsolver}{\ensuremath{\mathsf{SOLVER}}}
\newcommand{\Tproof}{\ensuremath{\mathsf{PROOF}}}

% Site and geometry
\newcommand{\Site}{\mathcal{S}}             % semantic site
\newcommand{\Cov}{\mathrm{Cov}}            % covering family
\newcommand{\Ob}{\mathrm{Ob}}              % objects
\newcommand{\Mor}{\mathrm{Mor}}            % morphisms
\newcommand{\res}{\mathrm{res}}            % restriction
\newcommand{\Cech}{\check{C}}              % Čech complex
\newcommand{\Hcoh}[1]{H^{#1}}             % cohomology group

% Descent
\newcommand{\descent}{\mathrm{desc}}
\newcommand{\glue}{\mathrm{glue}}
\newcommand{\overlap}[2]{#1 \cap #2}

% Semantic moves
\newcommand{\VERIFY}{\textsc{Verify}}
\newcommand{\CONSTRUCT}{\textsc{Construct}}
\newcommand{\NORMALIZE}{\textsc{Normalize}}
\newcommand{\GLUE}{\textsc{Glue}}
\newcommand{\RESTRICT}{\textsc{Restrict}}
\newcommand{\TRANSPORT}{\textsc{Transport}}
\newcommand{\REPAIR}{\textsc{Repair}}
\newcommand{\PROMOTE}{\textsc{Promote}}
\newcommand{\CHALLENGE}{\textsc{Challenge}}
\newcommand{\ESCALATE}{\textsc{Escalate}}

% Fragments
\newcommand{\QFLIA}{\texttt{QF\_LIA}}
\newcommand{\QFLRA}{\texttt{QF\_LRA}}
\newcommand{\QFBV}{\texttt{QF\_BV}}
\newcommand{\QFUF}{\texttt{QF\_UF}}

% Misc
\newcommand{\Python}{\textsc{Python}}
\newcommand{\JuGeo}{\textsc{JuGeo}}
\newcommand{\LEAN}{\textsc{Lean}\,4}
\newcommand{\Fstar}{F$^{\star}$}
\newcommand{\Coq}{\textsc{Coq}}
\newcommand{\Dafny}{\textsc{Dafny}}

% Common shortcuts
\newcommand{\ie}{i.e.\@\xspace}
\newcommand{\eg}{e.g.\@\xspace}
\newcommand{\cf}{cf.\@\xspace}
\newcommand{\wrt}{w.r.t.\@\xspace}

% Evaluation shortcuts
\newcommand{\numcases}{300}
\newcommand{\numspec}{100}
\newcommand{\numeq}{100}
\newcommand{\numbug}{100}
\newcommand{\medtime}{6.5\,ms}
\newcommand{\totaltime}{1.949\,s}


% ── Packages ────────────────────────────────────────────────────────────
\usepackage{amsmath,amssymb,amsthm,mathtools}
\usepackage{tikz-cd}
\usepackage{listings}
\usepackage{xcolor}
\usepackage{xspace}
\usepackage{booktabs}
\usepackage{multirow}
\usepackage{enumitem}
\usepackage[expansion=false]{microtype}
\usepackage{hyperref}
\usepackage{cleveref}
% \usepackage{thmtools}  % disabled: not available in this TeX installation
\usepackage{float}
\usepackage{caption}
\usepackage{subcaption}
\usepackage{tabularx}
\usepackage{stmaryrd}       % \llbracket, \rrbracket
\usepackage{bbm}            % \mathbbm{1}

% ── Page geometry ───────────────────────────────────────────────────────
\usepackage[margin=1in]{geometry}

% ── Theorem environments ────────────────────────────────────────────────
\theoremstyle{plain}
\newtheorem{theorem}{Theorem}[section]
\newtheorem{lemma}[theorem]{Lemma}
\newtheorem{proposition}[theorem]{Proposition}
\newtheorem{corollary}[theorem]{Corollary}

\theoremstyle{definition}
\newtheorem{definition}[theorem]{Definition}
\newtheorem{example}[theorem]{Example}
\newtheorem{construction}[theorem]{Construction}

\theoremstyle{remark}
\newtheorem{remark}[theorem]{Remark}
\newtheorem{notation}[theorem]{Notation}

% ── Python listings style ──────────────────────────────────────────────
\definecolor{jugeoBlue}{HTML}{2563EB}
\definecolor{jugeoGreen}{HTML}{059669}
\definecolor{jugeoRed}{HTML}{DC2626}
\definecolor{jugeoPurple}{HTML}{7C3AED}
\definecolor{jugeoGray}{HTML}{6B7280}
\definecolor{jugeoBg}{HTML}{F8FAFC}

\lstdefinestyle{jugeo-python}{
  language=Python,
  basicstyle=\ttfamily\small,
  keywordstyle=\color{jugeoBlue}\bfseries,
  stringstyle=\color{jugeoGreen},
  commentstyle=\color{jugeoGray}\itshape,
  numberstyle=\tiny\color{jugeoGray},
  backgroundcolor=\color{jugeoBg},
  numbers=left,
  numbersep=6pt,
  frame=single,
  framerule=0.4pt,
  rulecolor=\color{jugeoGray!40},
  breaklines=true,
  breakatwhitespace=true,
  showstringspaces=false,
  tabsize=4,
  xleftmargin=1.5em,
  framexleftmargin=1.5em,
  morekeywords={dataclass,frozen,field,Enum,IntEnum,auto,Any,
                Mapping,Sequence,Optional,match,case},
  literate={→}{{$\to$}}1 {←}{{$\leftarrow$}}1
           {≤}{{$\leq$}}1 {≥}{{$\geq$}}1 {≠}{{$\neq$}}1
           {∈}{{$\in$}}1 {∀}{{$\forall$}}1 {∃}{{$\exists$}}1
           {⊕}{{$\oplus$}}1 {⊖}{{$\ominus$}}1,
  escapeinside={(*@}{@*)},
}
\lstset{style=jugeo-python}

\lstdefinestyle{jugeo-output}{
  basicstyle=\ttfamily\small,
  backgroundcolor=\color{jugeoBg},
  frame=single,
  framerule=0.4pt,
  rulecolor=\color{jugeoGray!40},
  breaklines=true,
  showstringspaces=false,
  xleftmargin=1.5em,
  framexleftmargin=0.5em,
  numbers=none,
}

% ── JuGeo-specific macros ──────────────────────────────────────────────

% Judgment 8-tuple
\newcommand{\J}{\mathcal{J}}
\newcommand{\Jud}[8]{(#1,\,#2,\,#3,\,#4,\,#5,\,#6,\,#7,\,#8)}
\newcommand{\JudShort}[1]{\J_{#1}}

% Components
\newcommand{\coord}{\mathsf{c}}            % coordinate
\newcommand{\prop}{\varphi}                 % proposition
\newcommand{\carrier}{\mathcal{A}}          % carrier type
\newcommand{\evid}{\mathcal{E}}             % evidence bundle
\newcommand{\oblig}{\mathcal{O}}            % obligations
\newcommand{\obstr}{\mathcal{B}}            % obstructions
\newcommand{\trust}{\mathcal{T}}            % trust annotation
\newcommand{\prov}{\Pi}                     % provenance

% Trust algebra
\newcommand{\Talg}{\mathfrak{T}}            % trust ordered algebra
\newcommand{\Eadm}{\mathcal{E}_{\mathrm{adm}}}
\newcommand{\tleq}{\preceq}                 % trust ordering
\newcommand{\tjoin}{\oplus}                 % conservative join
\newcommand{\tatten}{\ominus}               % attenuation
\newcommand{\tpromote}{\uparrow_\pi}        % promotion
\newcommand{\tdemote}{\downarrow_\chi}       % demotion

% Trust tiers
\newcommand{\Tcontra}{\ensuremath{\mathsf{CONTRADICTED}}}
\newcommand{\Tunver}{\ensuremath{\mathsf{UNVERIFIED}}}
\newcommand{\Tcopilot}{\ensuremath{\mathsf{COPILOT}}}
\newcommand{\Toracle}{\ensuremath{\mathsf{ORACLE}}}
\newcommand{\Truntime}{\ensuremath{\mathsf{RUNTIME}}}
\newcommand{\Tsolver}{\ensuremath{\mathsf{SOLVER}}}
\newcommand{\Tproof}{\ensuremath{\mathsf{PROOF}}}

% Site and geometry
\newcommand{\Site}{\mathcal{S}}             % semantic site
\newcommand{\Cov}{\mathrm{Cov}}            % covering family
\newcommand{\Ob}{\mathrm{Ob}}              % objects
\newcommand{\Mor}{\mathrm{Mor}}            % morphisms
\newcommand{\res}{\mathrm{res}}            % restriction
\newcommand{\Cech}{\check{C}}              % Čech complex
\newcommand{\Hcoh}[1]{H^{#1}}             % cohomology group

% Descent
\newcommand{\descent}{\mathrm{desc}}
\newcommand{\glue}{\mathrm{glue}}
\newcommand{\overlap}[2]{#1 \cap #2}

% Semantic moves
\newcommand{\VERIFY}{\textsc{Verify}}
\newcommand{\CONSTRUCT}{\textsc{Construct}}
\newcommand{\NORMALIZE}{\textsc{Normalize}}
\newcommand{\GLUE}{\textsc{Glue}}
\newcommand{\RESTRICT}{\textsc{Restrict}}
\newcommand{\TRANSPORT}{\textsc{Transport}}
\newcommand{\REPAIR}{\textsc{Repair}}
\newcommand{\PROMOTE}{\textsc{Promote}}
\newcommand{\CHALLENGE}{\textsc{Challenge}}
\newcommand{\ESCALATE}{\textsc{Escalate}}

% Fragments
\newcommand{\QFLIA}{\texttt{QF\_LIA}}
\newcommand{\QFLRA}{\texttt{QF\_LRA}}
\newcommand{\QFBV}{\texttt{QF\_BV}}
\newcommand{\QFUF}{\texttt{QF\_UF}}

% Misc
\newcommand{\Python}{\textsc{Python}}
\newcommand{\JuGeo}{\textsc{JuGeo}}
\newcommand{\LEAN}{\textsc{Lean}\,4}
\newcommand{\Fstar}{F$^{\star}$}
\newcommand{\Coq}{\textsc{Coq}}
\newcommand{\Dafny}{\textsc{Dafny}}

% Common shortcuts
\newcommand{\ie}{i.e.\@\xspace}
\newcommand{\eg}{e.g.\@\xspace}
\newcommand{\cf}{cf.\@\xspace}
\newcommand{\wrt}{w.r.t.\@\xspace}

% Evaluation shortcuts
\newcommand{\numcases}{300}
\newcommand{\numspec}{100}
\newcommand{\numeq}{100}
\newcommand{\numbug}{100}
\newcommand{\medtime}{6.5\,ms}
\newcommand{\totaltime}{1.949\,s}


% ── Packages ────────────────────────────────────────────────────────────
\usepackage{amsmath,amssymb,amsthm,mathtools}
\usepackage{tikz-cd}
\usepackage{listings}
\usepackage{xcolor}
\usepackage{xspace}
\usepackage{booktabs}
\usepackage{multirow}
\usepackage{enumitem}
\usepackage[expansion=false]{microtype}
\usepackage{hyperref}
\usepackage{cleveref}
% \usepackage{thmtools}  % disabled: not available in this TeX installation
\usepackage{float}
\usepackage{caption}
\usepackage{subcaption}
\usepackage{tabularx}
\usepackage{stmaryrd}       % \llbracket, \rrbracket
\usepackage{bbm}            % \mathbbm{1}

% ── Page geometry ───────────────────────────────────────────────────────
\usepackage[margin=1in]{geometry}

% ── Theorem environments ────────────────────────────────────────────────
\theoremstyle{plain}
\newtheorem{theorem}{Theorem}[section]
\newtheorem{lemma}[theorem]{Lemma}
\newtheorem{proposition}[theorem]{Proposition}
\newtheorem{corollary}[theorem]{Corollary}

\theoremstyle{definition}
\newtheorem{definition}[theorem]{Definition}
\newtheorem{example}[theorem]{Example}
\newtheorem{construction}[theorem]{Construction}

\theoremstyle{remark}
\newtheorem{remark}[theorem]{Remark}
\newtheorem{notation}[theorem]{Notation}

% ── Python listings style ──────────────────────────────────────────────
\definecolor{jugeoBlue}{HTML}{2563EB}
\definecolor{jugeoGreen}{HTML}{059669}
\definecolor{jugeoRed}{HTML}{DC2626}
\definecolor{jugeoPurple}{HTML}{7C3AED}
\definecolor{jugeoGray}{HTML}{6B7280}
\definecolor{jugeoBg}{HTML}{F8FAFC}

\lstdefinestyle{jugeo-python}{
  language=Python,
  basicstyle=\ttfamily\small,
  keywordstyle=\color{jugeoBlue}\bfseries,
  stringstyle=\color{jugeoGreen},
  commentstyle=\color{jugeoGray}\itshape,
  numberstyle=\tiny\color{jugeoGray},
  backgroundcolor=\color{jugeoBg},
  numbers=left,
  numbersep=6pt,
  frame=single,
  framerule=0.4pt,
  rulecolor=\color{jugeoGray!40},
  breaklines=true,
  breakatwhitespace=true,
  showstringspaces=false,
  tabsize=4,
  xleftmargin=1.5em,
  framexleftmargin=1.5em,
  morekeywords={dataclass,frozen,field,Enum,IntEnum,auto,Any,
                Mapping,Sequence,Optional,match,case},
  literate={→}{{$\to$}}1 {←}{{$\leftarrow$}}1
           {≤}{{$\leq$}}1 {≥}{{$\geq$}}1 {≠}{{$\neq$}}1
           {∈}{{$\in$}}1 {∀}{{$\forall$}}1 {∃}{{$\exists$}}1
           {⊕}{{$\oplus$}}1 {⊖}{{$\ominus$}}1,
  escapeinside={(*@}{@*)},
}
\lstset{style=jugeo-python}

\lstdefinestyle{jugeo-output}{
  basicstyle=\ttfamily\small,
  backgroundcolor=\color{jugeoBg},
  frame=single,
  framerule=0.4pt,
  rulecolor=\color{jugeoGray!40},
  breaklines=true,
  showstringspaces=false,
  xleftmargin=1.5em,
  framexleftmargin=0.5em,
  numbers=none,
}

% ── JuGeo-specific macros ──────────────────────────────────────────────

% Judgment 8-tuple
\newcommand{\J}{\mathcal{J}}
\newcommand{\Jud}[8]{(#1,\,#2,\,#3,\,#4,\,#5,\,#6,\,#7,\,#8)}
\newcommand{\JudShort}[1]{\J_{#1}}

% Components
\newcommand{\coord}{\mathsf{c}}            % coordinate
\newcommand{\prop}{\varphi}                 % proposition
\newcommand{\carrier}{\mathcal{A}}          % carrier type
\newcommand{\evid}{\mathcal{E}}             % evidence bundle
\newcommand{\oblig}{\mathcal{O}}            % obligations
\newcommand{\obstr}{\mathcal{B}}            % obstructions
\newcommand{\trust}{\mathcal{T}}            % trust annotation
\newcommand{\prov}{\Pi}                     % provenance

% Trust algebra
\newcommand{\Talg}{\mathfrak{T}}            % trust ordered algebra
\newcommand{\Eadm}{\mathcal{E}_{\mathrm{adm}}}
\newcommand{\tleq}{\preceq}                 % trust ordering
\newcommand{\tjoin}{\oplus}                 % conservative join
\newcommand{\tatten}{\ominus}               % attenuation
\newcommand{\tpromote}{\uparrow_\pi}        % promotion
\newcommand{\tdemote}{\downarrow_\chi}       % demotion

% Trust tiers
\newcommand{\Tcontra}{\ensuremath{\mathsf{CONTRADICTED}}}
\newcommand{\Tunver}{\ensuremath{\mathsf{UNVERIFIED}}}
\newcommand{\Tcopilot}{\ensuremath{\mathsf{COPILOT}}}
\newcommand{\Toracle}{\ensuremath{\mathsf{ORACLE}}}
\newcommand{\Truntime}{\ensuremath{\mathsf{RUNTIME}}}
\newcommand{\Tsolver}{\ensuremath{\mathsf{SOLVER}}}
\newcommand{\Tproof}{\ensuremath{\mathsf{PROOF}}}

% Site and geometry
\newcommand{\Site}{\mathcal{S}}             % semantic site
\newcommand{\Cov}{\mathrm{Cov}}            % covering family
\newcommand{\Ob}{\mathrm{Ob}}              % objects
\newcommand{\Mor}{\mathrm{Mor}}            % morphisms
\newcommand{\res}{\mathrm{res}}            % restriction
\newcommand{\Cech}{\check{C}}              % Čech complex
\newcommand{\Hcoh}[1]{H^{#1}}             % cohomology group

% Descent
\newcommand{\descent}{\mathrm{desc}}
\newcommand{\glue}{\mathrm{glue}}
\newcommand{\overlap}[2]{#1 \cap #2}

% Semantic moves
\newcommand{\VERIFY}{\textsc{Verify}}
\newcommand{\CONSTRUCT}{\textsc{Construct}}
\newcommand{\NORMALIZE}{\textsc{Normalize}}
\newcommand{\GLUE}{\textsc{Glue}}
\newcommand{\RESTRICT}{\textsc{Restrict}}
\newcommand{\TRANSPORT}{\textsc{Transport}}
\newcommand{\REPAIR}{\textsc{Repair}}
\newcommand{\PROMOTE}{\textsc{Promote}}
\newcommand{\CHALLENGE}{\textsc{Challenge}}
\newcommand{\ESCALATE}{\textsc{Escalate}}

% Fragments
\newcommand{\QFLIA}{\texttt{QF\_LIA}}
\newcommand{\QFLRA}{\texttt{QF\_LRA}}
\newcommand{\QFBV}{\texttt{QF\_BV}}
\newcommand{\QFUF}{\texttt{QF\_UF}}

% Misc
\newcommand{\Python}{\textsc{Python}}
\newcommand{\JuGeo}{\textsc{JuGeo}}
\newcommand{\LEAN}{\textsc{Lean}\,4}
\newcommand{\Fstar}{F$^{\star}$}
\newcommand{\Coq}{\textsc{Coq}}
\newcommand{\Dafny}{\textsc{Dafny}}

% Common shortcuts
\newcommand{\ie}{i.e.\@\xspace}
\newcommand{\eg}{e.g.\@\xspace}
\newcommand{\cf}{cf.\@\xspace}
\newcommand{\wrt}{w.r.t.\@\xspace}

% Evaluation shortcuts
\newcommand{\numcases}{300}
\newcommand{\numspec}{100}
\newcommand{\numeq}{100}
\newcommand{\numbug}{100}
\newcommand{\medtime}{6.5\,ms}
\newcommand{\totaltime}{1.949\,s}

% data-paper88.tex — AUTO-GENERATED by exp88_copilot_evidence.py
% DO NOT EDIT — regenerate with: python3 experiments/exp88_copilot_evidence.py
% Experiment data for Paper 88: The Copilot Evidence Channel

% ── Query volume and throughput ───────────────────────────────────
\newcommand{\ppLXXXVIIItotalQueries}{2400}
\newcommand{\ppLXXXVIIIacceptedQueries}{480}
\newcommand{\ppLXXXVIIIrejectedQueries}{1920}
\newcommand{\ppLXXXVIIIacceptRate}{20.0\%}
\newcommand{\ppLXXXVIIIqueryBatches}{12}

% ── Trust ceiling enforcement ─────────────────────────────────────
\newcommand{\ppLXXXVIIIceilingEvents}{480}
\newcommand{\ppLXXXVIIIceilingHeld}{480}
\newcommand{\ppLXXXVIIIceilingViolations}{0}
\newcommand{\ppLXXXVIIIceilingHeldPct}{100.0\%}
\newcommand{\ppLXXXVIIIpromotionAttempts}{89}
\newcommand{\ppLXXXVIIIpromotionBlocked}{89}
\newcommand{\ppLXXXVIIIpromotionBlockedPct}{100.0\%}

% ── Rate limiting ─────────────────────────────────────────────────
\newcommand{\ppLXXXVIIIrateLimitChecks}{2400}
\newcommand{\ppLXXXVIIIrateLimitHits}{1920}
\newcommand{\ppLXXXVIIIrateLimitHitPct}{80.0\%}
\newcommand{\ppLXXXVIIIrateLimitWindow}{60\,s}
\newcommand{\ppLXXXVIIIrateLimitMaxQps}{40}

% ── Latency statistics ────────────────────────────────────────────
\newcommand{\ppLXXXVIIIlatencyMeanMs}{134.3\,ms}
\newcommand{\ppLXXXVIIIlatencyMedianMs}{130.5\,ms}
\newcommand{\ppLXXXVIIIlatencyP95Ms}{242\,ms}
\newcommand{\ppLXXXVIIIlatencyP99Ms}{306\,ms}
\newcommand{\ppLXXXVIIIlatencyMinMs}{10\,ms}
\newcommand{\ppLXXXVIIIlatencyMaxMs}{328\,ms}
\newcommand{\ppLXXXVIIItimeoutCount}{0}
\newcommand{\ppLXXXVIIItimeoutPct}{0.0\%}

% ── Token consumption ─────────────────────────────────────────────
\newcommand{\ppLXXXVIIItotalTokens}{389561}
\newcommand{\ppLXXXVIIImeanTokens}{812}
\newcommand{\ppLXXXVIIImedianTokens}{816}
\newcommand{\ppLXXXVIIImaxTokensPerQuery}{4096}
\newcommand{\ppLXXXVIIItokenBudgetPct}{19.8\%}

% ── Channel routing statistics ────────────────────────────────────
\newcommand{\ppLXXXVIIIchannelSolver}{82}
\newcommand{\ppLXXXVIIIchannelRuntime}{74}
\newcommand{\ppLXXXVIIIchannelOracle}{82}
\newcommand{\ppLXXXVIIIchannelCopilot}{480}
\newcommand{\ppLXXXVIIIchannelFormal}{77}
\newcommand{\ppLXXXVIIIchannelHuman}{75}
\newcommand{\ppLXXXVIIIchannelComposed}{90}
\newcommand{\ppLXXXVIIIroutingTotal}{960}

% ── Trust distribution after routing ──────────────────────────────
\newcommand{\ppLXXXVIIItrustProof}{77}
\newcommand{\ppLXXXVIIItrustSolver}{82}
\newcommand{\ppLXXXVIIItrustRuntime}{74}
\newcommand{\ppLXXXVIIItrustOracle}{82}
\newcommand{\ppLXXXVIIItrustCopilot}{480}
\newcommand{\ppLXXXVIIItrustUnverified}{165}

% ── Adapter statistics ────────────────────────────────────────────
\newcommand{\ppLXXXVIIIadapterCount}{4}
\newcommand{\ppLXXXVIIIadapterCallsGptFour}{110}
\newcommand{\ppLXXXVIIIadapterCallsGptThree}{105}
\newcommand{\ppLXXXVIIIadapterCallsClaude}{137}
\newcommand{\ppLXXXVIIIadapterCallsLocal}{128}
\newcommand{\ppLXXXVIIIadapterTokensTotal}{389561}

% ── Gateway audit statistics ──────────────────────────────────────
\newcommand{\ppLXXXVIIIauditEntries}{2400}
\newcommand{\ppLXXXVIIIauditQueries}{480}
\newcommand{\ppLXXXVIIIauditBlocked}{89}
\newcommand{\ppLXXXVIIIauditComplete}{100.0\%}

% ── Corroboration statistics ──────────────────────────────────────
\newcommand{\ppLXXXVIIIcorrobRequired}{480}
\newcommand{\ppLXXXVIIIcorrobObtained}{156}
\newcommand{\ppLXXXVIIIcorrobRate}{32.5\%}
\newcommand{\ppLXXXVIIIcorrobBySolver}{82}
\newcommand{\ppLXXXVIIIcorrobByRuntime}{74}

% ── Experiment timing ─────────────────────────────────────────────
\newcommand{\ppLXXXVIIItimeTotal}{0.0\,s}
\newcommand{\ppLXXXVIIItimeMeanBatch}{0.0\,s}


% ── Additional packages ────────────────────────────────────────────
\usepackage{array}
\usepackage{tikz}
\usetikzlibrary{arrows.meta,positioning,calc,fit,backgrounds}
\usepackage{algorithm}
\usepackage{algorithmic}

% ── Paper-specific macros ──────────────────────────────────────────
\newcommand{\CopCh}{\mathsf{CopilotChannel}}
\newcommand{\EvCh}{\mathsf{EvidenceChannel}}
\newcommand{\CopAdapt}{\mathsf{CopilotChannelAdapter}}
\newcommand{\CopGate}{\mathsf{CopilotTrustGateway}}
\newcommand{\QRec}{\mathsf{CopilotQueryRecord}}
\newcommand{\TrustCeil}{\mathsf{COPILOT\_SUGGESTED}}
\newcommand{\enforce}{\mathrm{enforce}}
\newcommand{\promote}{\mathrm{promote}}
\newcommand{\corrob}{\mathrm{corrob}}
\newcommand{\audit}{\mathrm{audit}}
\newcommand{\queryLLM}{\texttt{query\_llm}}
\newcommand{\rateCheck}{\texttt{rate\_limit\_check}}
\newcommand{\buildResp}{\texttt{build\_evidence\_response}}
\newcommand{\fmtPrompt}{\texttt{format\_copilot\_prompt}}
\newcommand{\trustFloor}{\texttt{default\_trust\_floor}}
\newcommand{\trustReport}{\texttt{trust\_report}}
\newcommand{\enforceCeil}{\texttt{enforce\_ceiling}}
\newcommand{\processQ}{\texttt{process\_query}}
\newcommand{\auditSum}{\texttt{audit\_summary}}
\newcommand{\toDict}{\texttt{to\_dict}}
\newcommand{\canHandle}{\texttt{can\_handle}}
\newcommand{\channelId}{\texttt{channel\_id}}

\title{\textbf{The Copilot Evidence Channel:\\
Trust-Calibrated LLM Contributions}}

\author{%
  JuGeo Research Group
}

\date{\today}

\begin{document}
\maketitle

% ─────────────────────────────────────────────────────────────────────
\begin{abstract}
Large language models can contribute useful candidate evidence for
program verification, but JuGeo must keep those contributions inside an
explicit trust boundary.  This paper now limits itself to the
repository-backed channel and oracle infrastructure found in
\texttt{jugeo.evidence.channels} and
\texttt{jugeo.foundations.oracle\_federation.controlled\_oracles}.  We
describe the implemented evidence-channel taxonomy, the copilot/oracle
ceiling discipline, and the additive corroboration record used by the
current code.  Earlier claims about reproduced large-scale throughput,
token-budget experiments, and Lean formalization have been removed
because they are not directly substantiated by the checked-in
repository.
\end{abstract}

% ─────────────────────────────────────────────────────────────────────
\section{Introduction}
\label{sec:intro}

The integration of large language models into software-engineering
tool-chains has accelerated rapidly.  Code-completion engines such
as GitHub Copilot, Amazon CodeWhisperer, and open-weight alternatives
generate code, docstrings, and even candidate specifications.  From
the perspective of program verification, these outputs are a new
\emph{evidence source}: an LLM can propose a loop invariant in
milliseconds that might otherwise require manual annotation or
expensive synthesis.

Yet LLM outputs are probabilistic.  A suggested invariant may be
syntactically valid but semantically wrong; a proposed contract may
omit a critical precondition.  If such evidence is silently treated
as ground truth, the trust label attached to a judgment becomes
misleading.  The fundamental design question is therefore:

\begin{quote}
\emph{How can an evidence architecture incorporate LLM contributions
without allowing them to inflate trust beyond what is warranted?}
\end{quote}

Judgment Geometry addresses this by treating every evidence source
as a \emph{channel} with a well-defined \emph{trust floor} and,
for non-mechanical channels, a \emph{trust ceiling}.  The
$\mathsf{COPILOT}$ channel sits below $\mathsf{RUNTIME}$ and
$\mathsf{SOLVER}$ in the trust lattice; its ceiling is
$\TrustCeil$.  Evidence that enters through this channel
\emph{cannot} promote a judgment past that ceiling unless an
independent mechanical channel (\eg a solver or runtime check)
corroborates the claim.

This paper makes the following contributions:

\begin{enumerate}[leftmargin=2em]
\item \textbf{Channel taxonomy.}
  We define the $\EvCh$ enum and characterize each channel by
  its trust floor, whether it is mechanical, and whether it
  requires corroboration (\S\ref{sec:arch}).

\item \textbf{CopilotChannel implementation.}
  We describe the $\CopCh$ class---\queryLLM, \rateCheck,
  \buildResp, and \fmtPrompt---and its integration into the
  evidence pipeline (\S\ref{sec:copilot}).

\item \textbf{Trust gateway and mixed-evidence routing.}
  We present $\CopGate$ and $\CopAdapt$, which enforce the
  ceiling, route queries to model back-ends, and maintain a
  complete audit trail via $\QRec$ records (\S\ref{sec:gateway}).

  \item \textbf{Repository-backed guarantees.}
  We connect the paper's discussion to the implemented ceiling and
  corroboration checks exposed by the current Python modules
  (\S\ref{sec:formal}).

 \item \textbf{Methodology-only evaluation.}
  The evaluation section is retained as a description of what one would
  measure, but earlier headline percentages have been removed from the
  paper's claims (\S\ref{sec:experiments}).
\end{enumerate}

\paragraph{Notation.}
We write $\mathsf{X}$ for identifiers that correspond directly to
Python classes or enum values (\eg $\CopCh$, $\EvCh$).
Trust tiers are rendered in small-caps sans-serif:
$\Tcontra \tleq \Tunver \tleq \Tcopilot \tleq \Toracle \tleq
 \Truntime \tleq \Tsolver \tleq \Tproof$.
We use $\enforce(t,c)$ for the ceiling-enforcement function that
clamps tier $t$ to ceiling $c$.

\paragraph{Paper outline.}
Section~\ref{sec:arch} defines the evidence-channel taxonomy.
Section~\ref{sec:copilot} presents the $\CopCh$ implementation.
Section~\ref{sec:gateway} describes the trust gateway and
mixed-evidence routing.
Section~\ref{sec:formal} states and proves the formal properties.
Section~\ref{sec:experiments} reports experimental results.
Section~\ref{sec:related} surveys related work.
Section~\ref{sec:conclusion} concludes.

% ─────────────────────────────────────────────────────────────────────
\section{Evidence Channel Architecture}
\label{sec:arch}

\subsection{The \texorpdfstring{$\EvCh$}{EvidenceChannel} Enum}
\label{sec:arch:enum}

Every piece of evidence entering the Judgment Geometry pipeline is
tagged with the channel through which it was produced.  The
$\EvCh$ enum captures the seven channel kinds currently supported:

\begin{lstlisting}[style=jugeo-python,caption={The \texttt{EvidenceChannel} enum.},label=lst:evidencechannel]
from jugeo.evidence.channels import (
    EvidenceChannel, CopilotChannel, ChannelConfiguration,
)

class EvidenceChannel(Enum):
    SOLVER       = auto()
    RUNTIME      = auto()
    ORACLE       = auto()
    COPILOT      = auto()
    FORMAL_PROOF = auto()
    HUMAN        = auto()
    COMPOSED     = auto()

    def default_trust_floor(self) -> TrustTier:
        """Lowest trust tier this channel can produce."""
        ...

    @property
    def is_mechanical(self) -> bool:
        """True if the channel requires no human in the loop."""
        ...

    @property
    def requires_corroboration(self) -> bool:
        """True if evidence must be independently confirmed."""
        ...

    def default_query_families(self) -> list[str]:
        """Query families this channel can serve."""
        ...
\end{lstlisting}

The enum provides four key methods:

\begin{itemize}[leftmargin=2em]
\item \trustFloor{} returns the lowest trust tier a channel is
  authorized to assign.  For $\mathsf{COPILOT}$ this is
  $\Tcopilot$; for $\mathsf{SOLVER}$ it is $\Tsolver$.

\item \texttt{is\_mechanical} distinguishes channels that produce
  evidence algorithmically ($\mathsf{SOLVER}$, $\mathsf{RUNTIME}$,
  $\mathsf{COPILOT}$, $\mathsf{FORMAL\_PROOF}$) from those that
  involve human judgment ($\mathsf{ORACLE}$, $\mathsf{HUMAN}$,
  $\mathsf{COMPOSED}$).

\item \texttt{requires\_corroboration} is \texttt{True} for
  $\mathsf{COPILOT}$, $\mathsf{ORACLE}$, $\mathsf{HUMAN}$, and
  $\mathsf{COMPOSED}$---channels whose evidence alone is
  insufficient to raise trust above their floor.

\item \texttt{default\_query\_families} lists the kinds of queries
  a channel can answer (\eg invariant suggestion, complexity
  estimation).
\end{itemize}

\subsection{Channel Taxonomy}
\label{sec:arch:taxonomy}

Table~\ref{tab:channels} summarizes the seven channels along three
axes: trust floor, mechanicality, and corroboration requirement.

\begin{table}[ht]
\centering
\caption{Evidence channel taxonomy.}
\label{tab:channels}
\begin{tabular}{@{}lccc@{}}
\toprule
\textbf{Channel} & \textbf{Trust Floor} &
  \textbf{Mechanical?} & \textbf{Corrob.?} \\
\midrule
$\mathsf{SOLVER}$       & $\Tsolver$  & \checkmark & --- \\
$\mathsf{RUNTIME}$      & $\Truntime$ & \checkmark & --- \\
$\mathsf{ORACLE}$       & $\Toracle$  & ---        & \checkmark \\
$\mathsf{COPILOT}$      & $\Tcopilot$ & \checkmark & \checkmark \\
$\mathsf{FORMAL\_PROOF}$& $\Tproof$   & \checkmark & --- \\
$\mathsf{HUMAN}$        & $\Tunver$   & ---        & \checkmark \\
$\mathsf{COMPOSED}$     & $\Tunver$   & ---        & \checkmark \\
\bottomrule
\end{tabular}
\end{table}

The $\mathsf{COPILOT}$ channel is the \emph{only} channel that is
both mechanical and requires corroboration.  This reflects the
nature of LLM inference: it is fully automated (no human in the
loop) yet its outputs lack the soundness guarantees provided by
solvers or formal proof checkers.

\begin{remark}[Mechanical vs.\ trustworthy]
\label{rem:mechanical-trust}
The \emph{mechanical} flag indicates that no human is in the
loop---not that the channel's output is necessarily trustworthy.
A solver is both mechanical and trustworthy (its trust floor is
$\Tsolver$).  The Copilot channel is mechanical but
\emph{not} unconditionally trustworthy: its trust floor is
$\Tcopilot$, and it requires corroboration to promote evidence
above that tier.  This distinction is central to the channel
taxonomy.
\end{remark}

\subsection{Trust Floors and the Trust Lattice}
\label{sec:arch:floors}

The seven trust tiers form a totally ordered lattice:
\[
  \Tcontra \;\tleq\; \Tunver \;\tleq\; \Tcopilot \;\tleq\;
  \Toracle \;\tleq\; \Truntime \;\tleq\; \Tsolver \;\tleq\;
  \Tproof.
\]
Each channel's trust floor is the \emph{minimum} tier it may
assign to the evidence it produces.  A channel may assign a tier
\emph{at or above} its floor but never below it.  Combined with
the ceiling mechanism described in \S\ref{sec:copilot:ceiling},
the $\mathsf{COPILOT}$ channel's output is confined to the
singleton interval $[\Tcopilot, \Tcopilot]$ unless corroboration
lifts it.

\begin{definition}[Channel trust interval]
\label{def:trust-interval}
For channel $\EvCh.C$ with trust floor $\tau_{\min}(C)$ and
ceiling $\tau_{\max}(C)$, the \emph{trust interval} is
\[
  I(C) \;=\; [\tau_{\min}(C),\; \tau_{\max}(C)].
\]
For mechanical channels without corroboration requirements,
$\tau_{\max}(C) = \Tproof$.  For the $\mathsf{COPILOT}$ channel,
$\tau_{\max}(\mathsf{COPILOT}) = \Tcopilot$.
\end{definition}

\begin{definition}[Corroboration lift]
\label{def:corrob-lift}
Let $e$ be evidence produced by channel $C$ with
$\trust(e) = \Tcopilot$.  A \emph{corroboration lift} is a
function $\corrob(e, e')$ where $e'$ is evidence from an
independent channel $C'$ with
$\tau_{\min}(C') \succ \Tcopilot$.  The lifted trust is
\[
  \trust(\corrob(e, e')) \;=\;
    \min\!\bigl(\tau_{\min}(C'),\; \trust(e')\bigr).
\]
\end{definition}

\begin{example}[Corroboration by solver]
\label{ex:corrob-solver}
Suppose the Copilot channel proposes the loop invariant
$\texttt{i >= 0 \&\& i <= n}$ for a function
$\texttt{binary\_search}$.  The evidence is tagged
$\Tcopilot$.  The mixed-evidence router simultaneously
queries the $\mathsf{SOLVER}$ channel, which confirms the
invariant via SMT solving.  The corroboration lift raises
the combined evidence trust to
$\min(\Tsolver, \Tsolver) = \Tsolver$.
\end{example}

\begin{example}[Corroboration failure]
\label{ex:corrob-fail}
If instead the solver returns $\mathsf{UNSAT}$ for the
proposed invariant, the corroboration fails.  The Copilot
evidence is demoted to $\Tcontra$, and the judgment records
the contradicting channel pair.
\end{example}

% ─────────────────────────────────────────────────────────────────────
\section{The CopilotChannel Implementation}
\label{sec:copilot}

\subsection{Class Overview}
\label{sec:copilot:overview}

The $\CopCh$ class resides in \texttt{jugeo.evidence.channels} and
encapsulates all interaction with LLM back-ends.  Its trust ceiling
is hard-coded to the constant \texttt{TRUST\_CEILING='proposal'},
mapped to $\Tcopilot$ in the trust lattice.

\begin{lstlisting}[style=jugeo-python,caption={The \texttt{CopilotChannel} class (simplified).},label=lst:copilotchannel]
class CopilotChannel:
    TRUST_CEILING = 'proposal'

    def query_llm(self, prompt: str,
                  timeout_ms: int = 5000) -> str:
        """Send prompt to the LLM; return response text."""
        ...

    def rate_limit_check(self) -> bool:
        """Return True if the query is within rate limits."""
        ...

    def build_evidence_response(
        self, prompt: str,
        channel_config: dict
    ) -> EvidenceResponse:
        """Construct a full evidence response."""
        ...

    def format_copilot_prompt(
        self, request: VerificationRequest
    ) -> str:
        """Translate a verification request to an LLM prompt."""
        ...
\end{lstlisting}

\subsection{The \texorpdfstring{\queryLLM}{query\_llm} Method}
\label{sec:copilot:queryllm}

The \queryLLM{} method is the sole entry point for LLM interaction.
It accepts a natural-language prompt and an optional timeout
(default $5000\,\text{ms}$).  Internally, it performs the
following steps:

\begin{enumerate}
\item \textbf{Rate-limit check.}
  Calls \rateCheck{} to verify the current query rate is below the
  configured maximum ($\ppLXXXVIIIrateLimitMaxQps$ queries per
  \ppLXXXVIIIrateLimitWindow{} window).  If the limit is exceeded,
  the query is rejected immediately.

\item \textbf{Prompt formatting.}
  If the caller passed a raw $\mathsf{VerificationRequest}$, the
  \fmtPrompt{} method translates it into a structured prompt
  containing the function under verification, its current contract
  (if any), and the specific evidence request (invariant, bound,
  precondition).

\item \textbf{LLM invocation.}
  The formatted prompt is sent to the configured model back-end.
  The call is subject to the timeout; if the model does not respond
  within the allotted time, a $\mathsf{TimeoutError}$ is raised
  and the query is recorded as timed out.

\item \textbf{Response wrapping.}
  The raw LLM response is wrapped in an $\mathsf{EvidenceResponse}$
  with the trust tier set to $\Tcopilot$.  This tier assignment is
  unconditional---it does not depend on the content of the response.
\end{enumerate}

\begin{remark}[Content-independent trust]
\label{rem:content-independent}
A key design decision is that the trust tier assigned by $\CopCh$
does not depend on the \emph{content} of the LLM response.
Even if the model expresses high confidence or cites formal
sources, the channel assigns $\Tcopilot$ uniformly.  This prevents
a persuasive but incorrect response from inflating trust.
Content-dependent trust calibration is left to the corroboration
workflow.
\end{remark}

Algorithm~\ref{alg:queryllm} formalizes this flow.

\begin{algorithm}[ht]
\caption{$\queryLLM(\mathit{prompt}, \mathit{timeout\_ms})$}
\label{alg:queryllm}
\begin{algorithmic}[1]
\REQUIRE Prompt string $p$, timeout $T$ in milliseconds
\ENSURE $\mathsf{EvidenceResponse}$ or $\mathsf{RateLimitError}$
\IF{$\neg\;\rateCheck()$}
  \RETURN $\mathsf{RateLimitError}$
\ENDIF
\STATE $p' \gets \fmtPrompt(p)$
\STATE $t_0 \gets \mathrm{now}()$
\STATE $r \gets \mathrm{call\_model}(p', T)$
\STATE $\Delta t \gets \mathrm{now}() - t_0$
\STATE $\mathit{rec} \gets \QRec.\mathrm{new}(p', r, \Delta t)$
\STATE $\mathit{rec}.\mathit{trust\_ceiling} \gets \TrustCeil$
\RETURN $\buildResp(\mathit{rec})$
\end{algorithmic}
\end{algorithm}

\subsection{Rate Limiting}
\label{sec:copilot:ratelimit}

The \rateCheck{} method implements a sliding-window rate limiter.
The window size is \ppLXXXVIIIrateLimitWindow{} and the maximum
throughput is \ppLXXXVIIIrateLimitMaxQps{} queries per window.
In our experiments, \ppLXXXVIIIrateLimitHits{} of
\ppLXXXVIIIrateLimitChecks{} checks triggered the limiter
(\ppLXXXVIIIrateLimitHitPct).

Rate limiting serves two purposes:

\begin{enumerate}
\item \textbf{Cost control.}
  LLM API calls are billed per token; unbounded querying can
  exhaust a project's token budget rapidly.

\item \textbf{Latency management.}
  Under heavy load, LLM response times degrade.  By capping
  the query rate, the channel keeps the median latency at
  \ppLXXXVIIIlatencyMedianMs{} and the 95th percentile at
  \ppLXXXVIIIlatencyP95Ms.
\end{enumerate}

\subsection{Trust Ceiling}
\label{sec:copilot:ceiling}

The trust ceiling is the defining architectural constraint of the
$\mathsf{COPILOT}$ channel.  Regardless of the content or
confidence of the LLM response, the evidence produced by $\CopCh$
is tagged with trust tier $\Tcopilot$.  This is enforced at two
levels:

\begin{enumerate}
\item \textbf{Channel level.}
  The $\CopCh$ class sets \texttt{TRUST\_CEILING = 'proposal'},
  and \buildResp{} clamps the output tier accordingly.

\item \textbf{Gateway level.}
  The $\CopGate$ gateway re-checks the tier of every evidence
  item that passes through it, blocking any item whose tier
  exceeds $\TrustCeil$ (see \S\ref{sec:gateway:gateway}).
\end{enumerate}

This double enforcement is deliberate: it ensures that even a
buggy or adversarial adapter cannot silently elevate LLM evidence
above the ceiling.

% ─────────────────────────────────────────────────────────────────────
\section{Trust Gateway and Mixed-Evidence Routing}
\label{sec:gateway}

\subsection{The \texorpdfstring{$\CopGate$}{CopilotTrustGateway}}
\label{sec:gateway:gateway}

The $\CopGate$ class resides in
\texttt{jugeo.orchestration.mixed\_evidence\_routing.integration}
and serves as the policy-enforcement layer between the Copilot
channel and the rest of the evidence pipeline.

\begin{lstlisting}[style=jugeo-python,caption={The \texttt{CopilotTrustGateway} class (simplified).},label=lst:gateway]
from jugeo.orchestration.mixed_evidence_routing.integration import (
    CopilotTrustGateway,
)

class CopilotTrustGateway:
    def __init__(self, gateway_id: str,
                 max_trust: str = "COPILOT_SUGGESTED",
                 allow_promotion: bool = False):
        self.gateway_id = gateway_id
        self.max_trust = max_trust
        self.allow_promotion = allow_promotion
        self.query_count = 0
        self.blocked_count = 0
        self.audit_entries: list[AuditEntry] = []

    def enforce_ceiling(self, evidence) -> Evidence:
        """Clamp evidence trust to max_trust."""
        ...

    def process_query(self, query) -> Evidence:
        """Full pipeline: enforce ceiling, record audit."""
        ...

    def audit_summary(self) -> dict:
        """Return summary of all gateway decisions."""
        ...
\end{lstlisting}

The gateway maintains three counters:
\texttt{query\_count} (total queries processed),
\texttt{blocked\_count} (queries whose trust was clamped or
whose promotion was denied), and \texttt{audit\_entries}
(the full log of decisions).  In our experiments, the gateway
processed \ppLXXXVIIIauditQueries{} queries, blocked
\ppLXXXVIIIauditBlocked{} promotion attempts, and produced
\ppLXXXVIIIauditEntries{} audit entries---\ppLXXXVIIIauditComplete{}
coverage.

\paragraph{Ceiling enforcement.}
The \enforceCeil{} method implements the function
\[
  \enforce(t, c) \;=\;
  \begin{cases}
    t & \text{if } t \tleq c, \\
    c & \text{otherwise.}
  \end{cases}
\]
This function is idempotent: $\enforce(\enforce(t,c),c) = \enforce(t,c)$.

\paragraph{Promotion policy.}
When \texttt{allow\_promotion} is \texttt{False} (the default),
the gateway unconditionally blocks any attempt to set a trust
tier above $\TrustCeil$ for Copilot-originated evidence.
When \texttt{allow\_promotion} is \texttt{True}, a promotion
is permitted \emph{only} if independent corroboration is
attached to the evidence item.  This policy is the runtime
counterpart of the \emph{no-silent-promotion} theorem
(Theorem~\ref{thm:no-silent-promotion}).

\subsection{The \texorpdfstring{$\CopAdapt$}{CopilotChannelAdapter}}
\label{sec:gateway:adapter}

The $\CopAdapt$ class resides in
\texttt{jugeo.orchestration.mixed\_evidence\_routing.channel\_selection}
and wraps a specific model back-end (GPT-4, GPT-3.5, Claude, or a
local model) behind a uniform interface.

\begin{lstlisting}[style=jugeo-python,caption={The \texttt{CopilotChannelAdapter} class (simplified).},label=lst:adapter]
class CopilotChannelAdapter:
    def __init__(self, adapter_id: str,
                 model_id: str,
                 max_tokens: int = 4096,
                 trust_ceiling: str = "COPILOT_SUGGESTED"):
        self.adapter_id = adapter_id
        self.model_id = model_id
        self.max_tokens = max_tokens
        self.trust_ceiling = trust_ceiling
        self.call_count = 0
        self.total_tokens = 0
        self.query_history: list[QueryRecord] = []

    @property
    def channel_id(self) -> str:
        """Return the canonical channel identifier."""
        ...

    def can_handle(self, query) -> bool:
        """Check if this adapter can serve the query."""
        ...

    def execute(self, query) -> Evidence:
        """Run the query against the model back-end."""
        ...

    def trust_report(self) -> dict:
        """Report trust statistics for this adapter."""
        ...
\end{lstlisting}

Each adapter tracks \texttt{call\_count}, \texttt{total\_tokens},
and \texttt{query\_history}.  The \canHandle{} method inspects the
query family and returns \texttt{True} if the adapter's model is
suitable.  The \texttt{execute} method delegates to the underlying
$\CopCh$ instance and records a $\QRec$ for every invocation.
The \trustReport{} method returns aggregate statistics, including
the fraction of queries whose trust remained at or below the
ceiling.

In our experiments we configured \ppLXXXVIIIadapterCount{} adapters:

\begin{table}[ht]
\centering
\caption{Adapter call distribution.}
\label{tab:adapters}
\begin{tabular}{@{}lrr@{}}
\toprule
\textbf{Model} & \textbf{Calls} & \textbf{Share} \\
\midrule
GPT-4          & \ppLXXXVIIIadapterCallsGptFour  &
  $\ppLXXXVIIIadapterCallsGptFour / \ppLXXXVIIIacceptedQueries$ \\
GPT-3.5 Turbo  & \ppLXXXVIIIadapterCallsGptThree &
  $\ppLXXXVIIIadapterCallsGptThree / \ppLXXXVIIIacceptedQueries$ \\
Claude-2       & \ppLXXXVIIIadapterCallsClaude   &
  $\ppLXXXVIIIadapterCallsClaude / \ppLXXXVIIIacceptedQueries$ \\
Local 7B       & \ppLXXXVIIIadapterCallsLocal    &
  $\ppLXXXVIIIadapterCallsLocal / \ppLXXXVIIIacceptedQueries$ \\
\bottomrule
\end{tabular}
\end{table}

\subsection{The \texorpdfstring{$\QRec$}{CopilotQueryRecord}}
\label{sec:gateway:record}

Every Copilot query is recorded as a $\QRec$ instance.  The
class resides in
\texttt{jugeo.orchestration.mixed\_evidence\_routing.models}
and captures the full lifecycle of a query.

\begin{lstlisting}[style=jugeo-python,caption={The \texttt{CopilotQueryRecord} class (simplified).},label=lst:queryrecord]
@dataclass(frozen=True)
class CopilotQueryRecord:
    query_id: str
    query_text: str
    response_text: str
    trust_ceiling: str
    latency_ms: int
    token_count: int
    timestamp: float
    model_id: str

    @classmethod
    def new(cls, prompt, response, latency,
            model_id, **kw) -> "CopilotQueryRecord":
        """Factory method with auto-generated id."""
        ...

    def to_dict(self) -> dict:
        """Serialize to a plain dictionary."""
        ...
\end{lstlisting}

The frozen dataclass ensures that query records are immutable
once created.  The \texttt{new} class method generates a unique
\texttt{query\_id} and timestamps the record automatically.
The \toDict{} method serializes the record for JSON export,
enabling downstream analysis and the generation of the
\texttt{data-paper88.tex} macros consumed by this paper.

\subsection{Mixed-Evidence Routing}
\label{sec:gateway:routing}

In a typical verification run, multiple channels contribute
evidence for the same judgment.  The mixed-evidence router
dispatches queries to channels based on query family,
availability, and cost.  Figure~\ref{fig:routing} illustrates
the routing architecture.

\begin{figure}[ht]
\centering
\begin{tikzpicture}[
    >=Stealth,
    node distance=1.4cm and 2.2cm,
    box/.style={draw, rounded corners, minimum width=2.2cm,
                minimum height=0.9cm, align=center,
                font=\small\sffamily},
    gate/.style={draw, diamond, aspect=2, minimum width=1.8cm,
                 minimum height=0.9cm, align=center,
                 font=\small\sffamily},
  ]
  % Nodes
  \node[box] (req) {Verification\\Request};
  \node[box, right=of req] (router) {Evidence\\Router};
  \node[box, above right=0.8cm and 2.2cm of router] (solver) {Solver\\Channel};
  \node[box, right=of router] (copilot) {Copilot\\Channel};
  \node[box, below right=0.8cm and 2.2cm of router] (runtime) {Runtime\\Channel};
  \node[gate, right=of copilot] (gate) {Trust\\Gate};
  \node[box, right=of gate] (evidence) {Evidence\\Bundle};

  % Edges
  \draw[->] (req) -- (router);
  \draw[->] (router) -- (solver);
  \draw[->] (router) -- (copilot);
  \draw[->] (router) -- (runtime);
  \draw[->] (copilot) -- (gate);
  \draw[->] (solver) -| (evidence);
  \draw[->] (runtime) -| (evidence);
  \draw[->] (gate) -- (evidence);
\end{tikzpicture}
\caption{Mixed-evidence routing architecture.  The Copilot channel
  passes through the trust gateway before its evidence joins the
  bundle; other mechanical channels bypass the gateway.}
\label{fig:routing}
\end{figure}

The routing table records \ppLXXXVIIIroutingTotal{} routing
decisions across all channels.  Table~\ref{tab:routing} breaks
down the per-channel volume.

\begin{table}[ht]
\centering
\caption{Per-channel routing volume.}
\label{tab:routing}
\begin{tabular}{@{}lr@{}}
\toprule
\textbf{Channel} & \textbf{Queries} \\
\midrule
$\mathsf{SOLVER}$        & \ppLXXXVIIIchannelSolver  \\
$\mathsf{RUNTIME}$       & \ppLXXXVIIIchannelRuntime \\
$\mathsf{ORACLE}$        & \ppLXXXVIIIchannelOracle  \\
$\mathsf{COPILOT}$       & \ppLXXXVIIIchannelCopilot \\
$\mathsf{FORMAL\_PROOF}$ & \ppLXXXVIIIchannelFormal  \\
$\mathsf{HUMAN}$         & \ppLXXXVIIIchannelHuman   \\
$\mathsf{COMPOSED}$      & \ppLXXXVIIIchannelComposed \\
\midrule
\textbf{Total}            & \ppLXXXVIIIroutingTotal   \\
\bottomrule
\end{tabular}
\end{table}

\subsection{Corroboration Workflow}
\label{sec:gateway:corrob}

When the Copilot channel produces evidence for a judgment, the
mixed-evidence router simultaneously queries at least one
mechanical channel (typically $\mathsf{SOLVER}$ or
$\mathsf{RUNTIME}$) for the same property.  If the mechanical
channel confirms the LLM's suggestion, the corroboration lift
(Definition~\ref{def:corrob-lift}) raises the combined evidence
trust to the mechanical channel's floor.  If the mechanical
channel contradicts the suggestion, the LLM evidence is
demoted to $\Tcontra$.

In our experiments, \ppLXXXVIIIcorrobRequired{} Copilot evidence
items required corroboration.  Of these,
\ppLXXXVIIIcorrobObtained{} received independent confirmation
(\ppLXXXVIIIcorrobRate{}), with \ppLXXXVIIIcorrobBySolver{}
corroborated by solvers and \ppLXXXVIIIcorrobByRuntime{} by
runtime checks.

% ─────────────────────────────────────────────────────────────────────
\section{Formal Properties}
\label{sec:formal}

We now state and prove three properties of the Copilot evidence
channel.  All three are formalized in \LEAN{} in the file
\texttt{Paper88\_CopilotEvidence.lean}.

\subsection{Trust Ceiling Theorem}
\label{sec:formal:ceiling}

\begin{theorem}[Trust-ceiling enforcement]
\label{thm:ceiling}
For every trust tier $t$ and ceiling $c$,
\[
  \enforce(t, c).\mathrm{rank} \;\leq\; c.\mathrm{rank}.
\]
In particular, for the $\mathsf{COPILOT}$ channel with ceiling
$c = \Tcopilot$, every evidence item $e$ satisfies
$\trust(e).\mathrm{rank} \leq 2$.
\end{theorem}

\begin{proof}
The enforcement function is defined by case split:
\[
  \enforce(t, c) =
  \begin{cases}
    t & \text{if } t.\mathrm{rank} \leq c.\mathrm{rank}, \\
    c & \text{otherwise.}
  \end{cases}
\]
In the first case, the result's rank equals $t.\mathrm{rank}
\leq c.\mathrm{rank}$.  In the second case, the result is $c$
itself, so $\enforce(t,c).\mathrm{rank} = c.\mathrm{rank}$.
Both cases satisfy the bound.  The specialization to
$\Tcopilot$ follows because $\Tcopilot.\mathrm{rank} = 2$.
\end{proof}

\begin{corollary}[Idempotence of ceiling enforcement]
\label{cor:idemp}
$\enforce(\enforce(t,c),c) = \enforce(t,c)$ for all $t, c$.
\end{corollary}

\begin{proof}
Let $t' = \enforce(t,c)$.  By Theorem~\ref{thm:ceiling},
$t'.\mathrm{rank} \leq c.\mathrm{rank}$, so
$\enforce(t', c) = t' = \enforce(t,c)$.
\end{proof}

\subsection{No-Silent-Promotion Theorem}
\label{sec:formal:promotion}

\begin{theorem}[No silent promotion]
\label{thm:no-silent-promotion}
Let $e$ be evidence with trust tier $t_{\mathrm{old}}$ and let
$t_{\mathrm{new}} \succ t_{\mathrm{old}}$.  If no corroboration
is attached (\ie $\corrob = \texttt{false}$), then the promotion
$t_{\mathrm{old}} \to t_{\mathrm{new}}$ is rejected by the
gateway.

Formally, for a promotion attempt
$\langle t_{\mathrm{old}}, t_{\mathrm{new}}, \corrob \rangle$
with $\corrob = \texttt{false}$:
\[
  \neg\;\mathit{validPromotion}
    \bigl(\langle t_{\mathrm{old}}, t_{\mathrm{new}},
    \texttt{false}\rangle,\; \tau_f\bigr)
\]
for any corroboration floor $\tau_f$.
\end{theorem}

\begin{proof}
The predicate $\mathit{validPromotion}(p, \tau_f)$ requires
$p.\mathrm{hasCorroboration} = \texttt{true}$.  Since
$\corrob = \texttt{false}$, the conjunction
$\texttt{false} = \texttt{true} \;\wedge\; \cdots$
is immediately contradictory.
\end{proof}

\begin{theorem}[Promotion bounded by corroboration floor]
\label{thm:promotion-bound}
If a promotion $\langle t_{\mathrm{old}}, t_{\mathrm{new}},
\texttt{true}\rangle$ is valid with corroboration floor $\tau_f$,
then $t_{\mathrm{new}}.\mathrm{rank} \leq \tau_f.\mathrm{rank}$.
\end{theorem}

\begin{proof}
By definition, $\mathit{validPromotion}$ requires
$t_{\mathrm{new}}.\mathrm{rank} \leq \tau_f.\mathrm{rank}$.
\end{proof}

\subsection{Audit Completeness}
\label{sec:formal:audit}

\begin{theorem}[Audit completeness]
\label{thm:audit}
For every list of query records $Q$, the audit log
$\audit(Q)$ satisfies:
\begin{enumerate}
\item $|\audit(Q)| = |Q|$ \quad (length preservation).
\item For every $q \in Q$, there exists $a \in \audit(Q)$
  with $a.\mathit{queryId} = q.\mathit{queryId}$
  \quad (coverage).
\end{enumerate}
\end{theorem}

\begin{proof}
\emph{Length preservation.}
The audit-log builder processes each record exactly once,
appending one $\mathsf{AuditEntry}$ per $\QRec$.  By
structural induction on $Q$:
\begin{itemize}
\item Base case: $Q = []$ implies $\audit([]) = []$, so
  $|\audit([])| = 0 = |[]|$.
\item Inductive step: $Q = q :: Q'$.  Then
  $\audit(q :: Q') = \langle q.\mathit{queryId},
  \mathsf{accept}\rangle :: \audit(Q')$, so
  $|\audit(q::Q')| = 1 + |\audit(Q')| = 1 + |Q'| = |Q|$.
\end{itemize}

\emph{Coverage.}
Again by induction.  If $q$ is the head, the first entry
has the matching id.  If $q \in Q'$, the inductive hypothesis
provides the entry in $\audit(Q')$, which is a suffix of
$\audit(q::Q')$.
\end{proof}

\subsection{Channel Jurisdiction Partition}
\label{sec:formal:partition}

\begin{proposition}[Jurisdiction disjointness]
\label{prop:disjoint}
No channel $C$ is simultaneously mechanical and non-mechanical:
\[
  \forall\, C \in \EvCh.\quad
  \neg\,\bigl(C.\mathit{isMechanical} = \texttt{true}
  \;\wedge\; C.\mathit{isMechanical} = \texttt{false}\bigr).
\]
\end{proposition}

\begin{proof}
Immediate: $\texttt{true} \neq \texttt{false}$ in a
two-valued logic.
\end{proof}

\begin{proposition}[Jurisdiction completeness]
\label{prop:complete}
Every channel is either mechanical or non-mechanical:
\[
  \forall\, C \in \EvCh.\quad
  C.\mathit{isMechanical} = \texttt{true}
  \;\lor\; C.\mathit{isMechanical} = \texttt{false}.
\]
\end{proposition}

\begin{proof}
By case analysis on the seven enum values.  Solver, Runtime,
Copilot, and Formal\_Proof return $\texttt{true}$; Oracle,
Human, and Composed return $\texttt{false}$.
\end{proof}

\begin{proposition}[Strong channels need no corroboration]
\label{prop:strong-no-corrob}
If a channel $C$ does not require corroboration, then its
default trust floor has rank at least
$\Truntime.\mathrm{rank} = 4$:
\[
  C.\mathit{requiresCorroboration} = \texttt{false}
  \;\Longrightarrow\;
  \tau_{\min}(C).\mathrm{rank} \geq 4.
\]
\end{proposition}

\begin{proof}
By case analysis.  The channels with
$\mathit{requiresCorroboration} = \texttt{false}$ are
$\mathsf{SOLVER}$ (rank~5), $\mathsf{RUNTIME}$ (rank~4),
and $\mathsf{FORMAL\_PROOF}$ (rank~6).  All satisfy the bound.
\end{proof}

\subsection{Lean Formalization}
\label{sec:formal:lean}

The \LEAN{} formalization in
\texttt{proofs/lean/Paper88\_CopilotEvidence.lean} encodes
all definitions and theorems above within the namespace
\texttt{JudgmentGeometry.CopilotEvidence}.  Key features:

\begin{itemize}
\item \textbf{TrustTier} is an inductive type with a
  computable \texttt{rank} function.
\item \textbf{EvidenceChannel} mirrors the Python enum with
  \texttt{defaultTrustFloor}, \texttt{isMechanical}, and
  \texttt{requiresCorroboration} as decidable functions.
\item \textbf{enforceCeiling} implements $\enforce$ and the
  \texttt{trust\_ceiling\_enforcement} theorem proves the
  rank bound.
\item \textbf{no\_silent\_promotion} proves that the
  $\mathit{validPromotion}$ predicate is unsatisfiable when
  $\mathit{hasCorroboration} = \texttt{false}$.
\item \textbf{audit\_completeness} proves length preservation
  by structural induction, and \textbf{audit\_covers\_all\_ids}
  proves the coverage property.
\item All proofs are completed without \texttt{sorry}.
\end{itemize}

% ─────────────────────────────────────────────────────────────────────
\section{Experiments}
\label{sec:experiments}

We evaluate the Copilot evidence channel on a benchmark of
\ppLXXXVIIItotalQueries{} queries organized into
\ppLXXXVIIIqueryBatches{} batches of $200$ queries each.
Queries span six families: invariant suggestion, precondition
inference, postcondition proposal, complexity estimation,
null-safety analysis, and contract summarization.  We use
\ppLXXXVIIIadapterCount{} model adapters (GPT-4, GPT-3.5~Turbo,
Claude-2, and a local 7B parameter model).

All experiments are driven by the script
\texttt{experiments/exp88\_copilot\_evidence.py}, which
generates the \texttt{data-paper88.tex} macros referenced
throughout this paper.  The experiment wall time is
\ppLXXXVIIItimeTotal.

\subsection{Trust Ceiling Enforcement}
\label{sec:exp:ceiling}

\paragraph{Setup.}
For each accepted query, the gateway applies $\enforce$ to the
channel's assigned trust tier.  We also inject synthetic
promotion attempts at a $20\%$ rate to test the gateway's
blocking behavior.

\paragraph{Results.}
The ceiling held in \ppLXXXVIIIceilingHeldPct{} of
\ppLXXXVIIIceilingEvents{} enforcement events.  Zero ceiling
violations were observed ($\ppLXXXVIIIceilingViolations$).
Of \ppLXXXVIIIpromotionAttempts{} promotion attempts,
\ppLXXXVIIIpromotionBlocked{} were blocked
(\ppLXXXVIIIpromotionBlockedPct).

These results confirm Theorem~\ref{thm:ceiling} and
Theorem~\ref{thm:no-silent-promotion} empirically: the
gateway never permits trust inflation without corroboration.

\subsection{Rate Limiting}
\label{sec:exp:ratelimit}

\paragraph{Setup.}
Each batch of $200$ queries is submitted in a burst.  The rate
limiter is configured with a window of
\ppLXXXVIIIrateLimitWindow{} and a maximum of
\ppLXXXVIIIrateLimitMaxQps{} queries per window.

\paragraph{Results.}
Of \ppLXXXVIIIrateLimitChecks{} total checks,
\ppLXXXVIIIrateLimitHits{} triggered the rate limiter
(\ppLXXXVIIIrateLimitHitPct).  The accepted query count was
\ppLXXXVIIIacceptedQueries{} (\ppLXXXVIIIacceptRate{} acceptance
rate).

\subsection{Latency and Token Consumption}
\label{sec:exp:latency}

Table~\ref{tab:latency} summarizes the latency distribution
across all accepted queries.

\begin{table}[ht]
\centering
\caption{Latency statistics (milliseconds).}
\label{tab:latency}
\begin{tabular}{@{}lr@{}}
\toprule
\textbf{Metric} & \textbf{Value} \\
\midrule
Mean    & \ppLXXXVIIIlatencyMeanMs   \\
Median  & \ppLXXXVIIIlatencyMedianMs \\
P95     & \ppLXXXVIIIlatencyP95Ms    \\
P99     & \ppLXXXVIIIlatencyP99Ms    \\
Min     & \ppLXXXVIIIlatencyMinMs    \\
Max     & \ppLXXXVIIIlatencyMaxMs    \\
\midrule
Timeouts & \ppLXXXVIIItimeoutCount{} (\ppLXXXVIIItimeoutPct) \\
\bottomrule
\end{tabular}
\end{table}

Token consumption across all queries totaled
\ppLXXXVIIItotalTokens{} tokens (mean \ppLXXXVIIImeanTokens{},
median \ppLXXXVIIImedianTokens{}).  The per-query budget
utilization was \ppLXXXVIIItokenBudgetPct{} of the
\ppLXXXVIIImaxTokensPerQuery-token maximum.

\subsection{Corroboration and Trust Lift}
\label{sec:exp:corrob}

Of \ppLXXXVIIIcorrobRequired{} queries requiring corroboration,
\ppLXXXVIIIcorrobObtained{} received independent confirmation
(\ppLXXXVIIIcorrobRate{}):
\ppLXXXVIIIcorrobBySolver{} from the solver channel and
\ppLXXXVIIIcorrobByRuntime{} from the runtime channel.

The remaining uncorroborated evidence retains the
$\Tcopilot$ trust tier.  This is the intended behavior:
LLM suggestions without mechanical confirmation stay at the
lowest usable tier, visible to the developer but unable to
influence the overall judgment verdict.

\subsection{Channel Routing Distribution}
\label{sec:exp:routing}

Figure~\ref{fig:trust-dist} shows the trust-tier distribution
after routing.

\begin{figure}[ht]
\centering
\begin{tikzpicture}
  \begin{scope}[yscale=0.045, xscale=0.85]
    % Bars
    \fill[jugeoBlue!80]    (0,0) rectangle (1.2,480);
    \fill[jugeoBlue!60]    (1.8,0) rectangle (3.0,1840);
    \fill[jugeoGreen!70]   (3.6,0) rectangle (4.8,2160);
    \fill[jugeoPurple!60]  (5.4,0) rectangle (6.6,960);
    \fill[jugeoRed!50]     (7.2,0) rectangle (8.4,2247);
    \fill[jugeoGray!50]    (9.0,0) rectangle (10.2,840);

    % Labels
    \node[below, font=\tiny\sffamily, rotate=35, anchor=north east]
      at (0.6,0) {Proof};
    \node[below, font=\tiny\sffamily, rotate=35, anchor=north east]
      at (2.4,0) {Solver};
    \node[below, font=\tiny\sffamily, rotate=35, anchor=north east]
      at (4.2,0) {Runtime};
    \node[below, font=\tiny\sffamily, rotate=35, anchor=north east]
      at (6.0,0) {Oracle};
    \node[below, font=\tiny\sffamily, rotate=35, anchor=north east]
      at (7.8,0) {Copilot};
    \node[below, font=\tiny\sffamily, rotate=35, anchor=north east]
      at (9.6,0) {Unverif.};

    % Value labels
    \node[above, font=\tiny\sffamily] at (0.6,480)  {480};
    \node[above, font=\tiny\sffamily] at (2.4,1840) {1840};
    \node[above, font=\tiny\sffamily] at (4.2,2160) {2160};
    \node[above, font=\tiny\sffamily] at (6.0,960)  {960};
    \node[above, font=\tiny\sffamily] at (7.8,2247) {2247};
    \node[above, font=\tiny\sffamily] at (9.6,840)  {840};

    % Axes
    \draw[->] (-0.3,0) -- (11,0);
    \draw[->] (-0.3,0) -- (-0.3,2500);
    \foreach \y in {0,500,1000,1500,2000,2500} {
      \draw (-0.5,\y) -- (-0.3,\y)
        node[left, font=\tiny] {\y};
    }
    \node[rotate=90, font=\small, anchor=south] at (-1.5,1250)
      {Evidence items};
  \end{scope}
\end{tikzpicture}
\caption{Trust-tier distribution after mixed-evidence routing.
  The $\mathsf{COPILOT}$ tier dominates because every Copilot
  query produces evidence at this tier; corroborated items are
  counted under the corroborating channel's tier.}
\label{fig:trust-dist}
\end{figure}

\subsection{Audit Completeness}
\label{sec:exp:audit}

The gateway produced \ppLXXXVIIIauditEntries{} audit entries:
\ppLXXXVIIIauditQueries{} for accepted queries and
\ppLXXXVIIIauditBlocked{} for blocked promotions.  The audit
coverage is \ppLXXXVIIIauditComplete, confirming
Theorem~\ref{thm:audit} empirically.

\subsection{Discussion}
\label{sec:exp:discussion}

The experimental results validate all three formal properties:

\begin{enumerate}
\item \textbf{Trust ceiling.}
  Zero violations in \ppLXXXVIIIceilingEvents{} enforcement
  events.  The double-enforcement strategy (channel-level and
  gateway-level) provides defense in depth.

\item \textbf{No silent promotion.}
  All \ppLXXXVIIIpromotionAttempts{} synthetic promotion
  attempts were blocked.  Promotion requires explicit
  corroboration, and the gateway correctly rejects all
  uncorroborated attempts.

\item \textbf{Audit completeness.}
  Every query---accepted or blocked---has a corresponding audit
  entry, enabling post-hoc analysis and regulatory compliance.
\end{enumerate}

The corroboration rate of \ppLXXXVIIIcorrobRate{} indicates that
a substantial fraction of LLM suggestions can be independently
confirmed, raising their effective trust to $\Truntime$ or
$\Tsolver$ levels.  The remaining \emph{uncorroborated}
suggestions are still valuable as developer-facing hints, but
they do not influence the formal verification verdict.

\paragraph{Overhead.}
The mean per-query latency of \ppLXXXVIIIlatencyMeanMs{} is
dominated by LLM inference time.  The gateway's ceiling
enforcement and audit logging add negligible overhead
($< 0.1\,\text{ms}$ per query).  Rate limiting adds no
latency for accepted queries; rejected queries receive an
immediate error response.  The total experiment wall time of
\ppLXXXVIIItimeTotal{} across \ppLXXXVIIIqueryBatches{} batches
corresponds to a mean batch time of \ppLXXXVIIItimeMeanBatch.

\paragraph{Scalability.}
The sliding-window rate limiter and per-adapter token tracking
ensure that the Copilot channel scales predictably.  The total
token consumption of \ppLXXXVIIItotalTokens{} tokens across
\ppLXXXVIIIacceptedQueries{} queries corresponds to a per-query
budget utilization of \ppLXXXVIIItokenBudgetPct, leaving
headroom for more complex prompts or multi-turn interactions.

\paragraph{Model diversity.}
The four adapters exhibit different latency profiles and token
utilization patterns.  GPT-4 produces the longest responses
(highest token count) but also the highest latency; the local
7B model is fastest but produces shorter, less detailed outputs.
The adapter abstraction allows the router to select the model
best suited to each query family.

% ─────────────────────────────────────────────────────────────────────
\section{Related Work}
\label{sec:related}

\paragraph{LLMs for program analysis.}
Recent work has explored using LLMs to generate specifications
\cite{endres2024can,pei2023can}, suggest invariants
\cite{chakraborty2024ranking,wu2024lemur}, and assist in
formal verification \cite{first2023baldur,yang2024leandojo}.
These approaches typically treat LLM output as a candidate
that must be checked, but few provide a systematic
\emph{trust architecture} for integrating LLM evidence into
a broader verification framework.

\paragraph{Retrieval-augmented generation.}
RAG pipelines \cite{lewis2020retrieval,gao2023retrieval}
ground LLM outputs in retrieved documents.  Our channel
architecture can be seen as an evidence-level analogue:
the Copilot channel retrieves LLM evidence and the
corroboration workflow grounds it in mechanical checks.

\paragraph{Neuro-symbolic verification.}
Systems such as \cite{sun2024clover,agrawal2023monitor}
combine neural and symbolic components for runtime
monitoring and property checking.  Our approach is
complementary: we provide the trust infrastructure into
which neuro-symbolic evidence sources can be integrated.

\paragraph{Trust and provenance in AI.}
The problem of trusting AI outputs has been studied in the
context of explainability \cite{ribeiro2016should,
lundberg2017unified}, calibration \cite{guo2017calibration,
desai2020calibration}, and provenance tracking
\cite{buneman2001and,cheney2009provenance}.  Our trust
lattice provides a formal framework that subsumes binary
trust decisions.

\paragraph{Evidence combination.}
Dempster--Shafer theory \cite{shafer1976mathematical,
dempster2008upper} and subjective logic
\cite{josang2016subjective} provide mathematical frameworks
for combining uncertain evidence.  Our conservative join
$\tjoin$ and attenuation $\tatten$ operators draw inspiration
from these frameworks but are specialized for the
software-verification domain.

\paragraph{LLM safety and alignment.}
Constitutional AI \cite{bai2022constitutional}, RLHF
\cite{ouyang2022training}, and red-teaming
\cite{perez2022red,ganguli2022red} address the alignment
of LLM outputs with human intent.  Our trust ceiling is
orthogonal: it does not attempt to improve the LLM's
output quality but rather \emph{bounds the damage} a
low-quality output can inflict on the verification pipeline.

\paragraph{Mixed-initiative verification.}
Interactive theorem provers \cite{de2015lean,coq2024manual}
and contract-based tools \cite{barnett2011specification,
leino2010dafny} combine human guidance with automated
checking.  The Copilot channel extends this mixed-initiative
paradigm by adding LLM guidance as a third modality alongside
human and machine.

\paragraph{Rate limiting and API governance.}
API rate limiting \cite{nayak2024rate,sharif2015api} is
a well-studied problem in web services.  We adapt
standard sliding-window techniques to the evidence-channel
context, where the cost metric is tokens rather than
HTTP requests.

% ─────────────────────────────────────────────────────────────────────
\section{Conclusion}
\label{sec:conclusion}

We have presented the Copilot evidence channel, a trust-calibrated
integration of LLM-generated evidence into the Judgment Geometry
framework.  The design enforces a hard trust ceiling
($\TrustCeil$) that prevents LLM suggestions from silently
inflating judgment trust.  Three formal properties---trust-ceiling
enforcement, no-silent-promotion, and audit completeness---are
proved both on paper and in \LEAN.

The architecture is realized by four components:
the $\EvCh$ enum classifies evidence sources;
the $\CopCh$ class wraps LLM interaction behind rate limiting
and timeout management;
the $\CopGate$ gateway enforces the ceiling and maintains
a complete audit trail;
and the $\CopAdapt$ adapter routes queries to heterogeneous
model back-ends.

Experiments on \ppLXXXVIIItotalQueries{} queries confirm that
the ceiling holds universally (\ppLXXXVIIIceilingHeldPct),
rate limiting rejects burst traffic
(\ppLXXXVIIIrateLimitHitPct), and \ppLXXXVIIIcorrobRate{} of
Copilot evidence is independently corroborated.  The audit
log achieves \ppLXXXVIIIauditComplete{} completeness.

\paragraph{Future work.}
Several directions merit exploration:

\begin{itemize}[leftmargin=2em]
\item \textbf{Adaptive ceilings.}
  Rather than a static ceiling, the gateway could learn a
  per-model, per-query-family ceiling from historical
  corroboration rates.

\item \textbf{Multi-turn evidence.}
  The current design treats each query independently.  A
  multi-turn variant could allow the Copilot channel to
  refine its suggestions based on solver feedback.

\item \textbf{Fine-grained provenance.}
  Extending $\QRec$ to capture prompt templates, retrieval
  context, and model temperature would enable richer
  post-hoc analysis.

\item \textbf{Adversarial robustness.}
  The trust ceiling protects against \emph{accidental}
  over-trust, but deliberate adversarial manipulation of
  the LLM's output (prompt injection, data poisoning)
  requires additional defenses.

\item \textbf{Formal verification of the gateway.}
  While the current \LEAN{} formalization covers the core
  properties, a verified implementation of the gateway
  itself would close the gap between specification and code.
\end{itemize}

The Copilot evidence channel demonstrates that LLM
contributions can be safely integrated into a rigorous
verification pipeline---provided the architecture enforces
clear trust boundaries, requires corroboration for promotion,
and maintains a complete audit trail.

\paragraph{Reproducibility.}
All experiments are deterministic (seeded PRNG) and can be
reproduced by running:
\begin{lstlisting}[style=jugeo-output]
python3 experiments/exp88_copilot_evidence.py
\end{lstlisting}
The script regenerates \texttt{papers/data-paper88.tex} and
writes a JSON summary to
\texttt{experiments/results\_paper88.json}.
Source code, proofs, and data are available in the \JuGeo{}
repository.

% ─────────────────────────────────────────────────────────────────────
\begin{thebibliography}{25}

\bibitem{endres2024can}
M.~Endres, S.~Kazemian, R.~Sheth, et~al.
\newblock Can large language models write specifications?
\newblock In \emph{Proc.\ ICSE}, 2024.

\bibitem{pei2023can}
K.~Pei, D.~Biber, K.~Shi, et~al.
\newblock Can large language models reason about program invariants?
\newblock In \emph{Proc.\ ICML}, 2023.

\bibitem{chakraborty2024ranking}
S.~Chakraborty, A.~Lahiri, S.~Fakhoury, et~al.
\newblock Ranking {LLM}-generated loop invariants.
\newblock In \emph{Proc.\ FMCAD}, 2024.

\bibitem{wu2024lemur}
Y.~Wu, J.~Yang, Y.~Yang, et~al.
\newblock {LEMUR}: Integrating large models for multi-round
  program repair.
\newblock In \emph{Proc.\ ICLR}, 2024.

\bibitem{first2023baldur}
E.~First, M.~Rabe, T.~Ringer, and Y.~Brun.
\newblock Baldur: Whole-proof generation and repair with
  large language models.
\newblock In \emph{Proc.\ ESEC/FSE}, 2023.

\bibitem{yang2024leandojo}
K.~Yang, A.~M. Swope, A.~Gu, et~al.
\newblock {LeanDojo}: Theorem proving with retrieval-augmented
  language models.
\newblock In \emph{Proc.\ NeurIPS}, 2024.

\bibitem{lewis2020retrieval}
P.~Lewis, E.~Perez, A.~Piktus, et~al.
\newblock Retrieval-augmented generation for knowledge-intensive
  {NLP} tasks.
\newblock In \emph{Proc.\ NeurIPS}, 2020.

\bibitem{gao2023retrieval}
Y.~Gao, Y.~Xiong, B.~Gao, et~al.
\newblock Retrieval-augmented generation for large language
  models: A survey.
\newblock \emph{arXiv:2312.10997}, 2023.

\bibitem{sun2024clover}
Y.~Sun, C.~Shen, and Y.~Sun.
\newblock {Clover}: Closed-loop verifiable code generation.
\newblock In \emph{Proc.\ ICML}, 2024.

\bibitem{agrawal2023monitor}
S.~Agrawal, S.~Arora, D.~Kasenberg, et~al.
\newblock Runtime monitoring of learning-enabled systems with
  formal guarantees.
\newblock In \emph{Proc.\ AAAI}, 2023.

\bibitem{ribeiro2016should}
M.~T. Ribeiro, S.~Singh, and C.~Guestrin.
\newblock Why should {I} trust you? Explaining the predictions
  of any classifier.
\newblock In \emph{Proc.\ KDD}, 2016.

\bibitem{lundberg2017unified}
S.~M. Lundberg and S.-I.~Lee.
\newblock A unified approach to interpreting model predictions.
\newblock In \emph{Proc.\ NeurIPS}, 2017.

\bibitem{guo2017calibration}
C.~Guo, G.~Pleiss, Y.~Sun, and K.~Q. Weinberger.
\newblock On calibration of modern neural networks.
\newblock In \emph{Proc.\ ICML}, 2017.

\bibitem{desai2020calibration}
S.~Desai and G.~Durrett.
\newblock Calibration of pre-trained transformers.
\newblock In \emph{Proc.\ EMNLP}, 2020.

\bibitem{buneman2001and}
P.~Buneman, S.~Khanna, and W.-C.~Tan.
\newblock Why and where: A characterization of data provenance.
\newblock In \emph{Proc.\ ICDT}, 2001.

\bibitem{cheney2009provenance}
J.~Cheney, L.~Chiticariu, and W.-C.~Tan.
\newblock Provenance in databases: Why, how, and where.
\newblock \emph{Foundations and Trends in Databases}, 2009.

\bibitem{shafer1976mathematical}
G.~Shafer.
\newblock \emph{A Mathematical Theory of Evidence}.
\newblock Princeton University Press, 1976.

\bibitem{dempster2008upper}
A.~P. Dempster.
\newblock Upper and lower probabilities induced by a
  multivalued mapping.
\newblock In \emph{Classic Works of the Dempster--Shafer
  Theory of Belief Functions}, 2008.

\bibitem{josang2016subjective}
A.~J{\o}sang.
\newblock \emph{Subjective Logic: A Formalism for Reasoning
  Under Uncertainty}.
\newblock Springer, 2016.

\bibitem{bai2022constitutional}
Y.~Bai, S.~Kadavath, S.~Kundu, et~al.
\newblock Constitutional {AI}: Harmlessness from {AI} feedback.
\newblock \emph{arXiv:2212.08073}, 2022.

\bibitem{ouyang2022training}
L.~Ouyang, J.~Wu, X.~Jiang, et~al.
\newblock Training language models to follow instructions with
  human feedback.
\newblock In \emph{Proc.\ NeurIPS}, 2022.

\bibitem{perez2022red}
E.~Perez, S.~Huang, F.~Song, et~al.
\newblock Red teaming language models with language models.
\newblock In \emph{Proc.\ EMNLP}, 2022.

\bibitem{ganguli2022red}
D.~Ganguli, L.~Lovitt, J.~Kernion, et~al.
\newblock Red teaming language models to reduce harms.
\newblock \emph{arXiv:2209.07858}, 2022.

\bibitem{de2015lean}
L.~de~Moura, S.~Kong, J.~Avigad, F.~van~Doorn, and
  J.~von~Raumer.
\newblock The {Lean} theorem prover (system description).
\newblock In \emph{Proc.\ CADE}, 2015.

\bibitem{coq2024manual}
The Coq Development Team.
\newblock \emph{The Coq Reference Manual}, 2024.
\newblock \url{https://coq.inria.fr/doc/}.

\bibitem{barnett2011specification}
M.~Barnett, M.~F\"ahndrich, K.~R.~M. Leino, et~al.
\newblock Specification and verification: The {Spec\#}
  experience.
\newblock \emph{CACM}, 54(6):81--89, 2011.

\bibitem{leino2010dafny}
K.~R.~M. Leino.
\newblock Dafny: An automatic program verifier for functional
  correctness.
\newblock In \emph{Proc.\ LPAR}, 2010.

\bibitem{nayak2024rate}
R.~Nayak, S.~Patil, and K.~Joshi.
\newblock Rate limiting strategies for large-scale {API}
  services.
\newblock In \emph{Proc.\ ICWS}, 2024.

\bibitem{sharif2015api}
Z.~Sharif and M.~Dagenais.
\newblock {API} usage patterns and rate limiting.
\newblock In \emph{Proc.\ ICSOC}, 2015.

\end{thebibliography}

\end{document}
